Özet

(TR)


Bu proje, güncel haber içeriklerinin çok dilli özetleme modelleriyle otomatik olarak işlenmesini, farklı model çıktılarının ölçülebilir kriterlerle değerlendirilmesini ve seçilen modellerin anlık haber akışı sağlayan bir web uygulamasına entegre edilmesini konu almaktadır. Projenin temel çıkış noktası, günümüzde haber içeriklerinin çok hızlı üretilmesi ve kullanıcıların farklı dillerdeki haberleri kısa, anlaşılır ve güvenilir biçimde takip etme ihtiyacının artmasıdır. Haber siteleri ve RSS akışları sürekli güncellenmekte; ancak bu akıştaki metinler çoğu zaman uzun, tekrar eden ayrıntılar içeren ve kullanıcı açısından hızlı taramaya elverişli olmayan yapıdadır. Bu nedenle proje, haber metinlerinden otomatik özet üretme problemini yalnızca bir model eğitimi problemi olarak değil, uçtan uca çalışan bir haber toplama, özetleme, değerlendirme ve sunma sistemi olarak ele almıştır.

Projenin kapsamı üç ana bileşenden oluşmaktadır. İlk bileşen, haber özetleme modellerinin eğitilmesi ve kullanılabilir hale getirilmesidir. Bu kapsamda BART, mBART ve mT5 tabanlı özetleme modelleri incelenmiş; özellikle çok dilli haber özetleme görevine uygun modeller üzerinde durulmuştur. Projede İngilizce odaklı modellerin yanında Türkçe, İngilizce, Fransızca, İspanyolca, Arapça, Hintçe, Çince, Japonca, Korece, Rusça ve Vietnamca gibi farklı dillerde çalışabilen çok dilli modeller kullanılmıştır. İkinci bileşen, eğitilen veya projeye dahil edilen modellerin sistematik biçimde değerlendirilmesi ve karşılaştırılmasıdır. Bu aşamada modeller yalnızca klasik metin benzerliği metrikleriyle değil, aynı zamanda özetin kaynak metne bağlılığı, tekrar oranı, okunabilirliği, sıkıştırma başarısı ve çalışma süresi gibi pratik kullanım açısından önemli ölçütlerle de incelenmiştir. Üçüncü bileşen ise bu modellerin anlık haber akışı sağlayan bir web uygulamasına entegre edilmesidir. Uygulama, RSS kaynaklarından haberleri toplamakta, haber sayfasından ana metni çıkarmakta, seçilen model ile özet üretmekte, sonuçları veritabanına kaydetmekte ve kullanıcıya filtrelenebilir bir arayüz üzerinden sunmaktadır.

Projenin sınırları da bu kapsamla uyumlu olarak belirlenmiştir. Çalışma, haber metinlerinin otomatik özetlenmesine odaklanmakta; haberlerin doğruluğunu bağımsız biçimde teyit eden bir fact-checking sistemi geliştirmemektedir. Üretilen özetlerin kaynak metne bağlı kalması için çeşitli temizlik, son işleme ve değerlendirme adımları uygulanmış olsa da sistem, haberin dış dünyadaki olgusal doğruluğunu garanti etmeyi amaçlamamaktadır. Benzer şekilde proje, bütün interneti tarayan genel amaçlı bir arama motoru değildir; önceden tanımlanmış ve doğrulanmış haber kaynaklarından gelen RSS akışları üzerinden çalışmaktadır. Haber sayfasından metin çıkarma işlemi, kaynak sitenin teknik yapısına, erişilebilirliğine ve sayfa şablonuna bağlıdır. Bu nedenle bazı kaynaklarda tam haber gövdesi alınamayabilir veya kaynak kalitesi düşükse sistem ilgili içeriği işlem dışında bırakabilir. Ayrıca çalışma, büyük ölçekli üretim ortamı optimizasyonundan ziyade araştırma ve prototip düzeyinde çalışan, ölçülebilir ve genişletilebilir bir sistem geliştirmeye odaklanmıştır.

Bu projenin temel araştırma problemi şu şekilde ifade edilebilir: Çok dilli haber metinleri için eğitilen veya uyarlanan özetleme modelleri, kalite, kaynak sadakati, hız ve kullanılabilirlik açısından nasıl karşılaştırılabilir ve bu modeller gerçek zamanlı haber akışı sunan bir web uygulamasında nasıl işlevsel biçimde kullanılabilir? Bu problem iki alt soruyu da beraberinde getirmektedir. Birincisi, farklı model mimarileri ve eğitim stratejileri haber özetleme kalitesini nasıl etkilemektedir? İkincisi, model çıktılarının yalnızca akademik değerlendirme ortamında değil, gerçek haber kaynaklarından gelen değişken ve gürültülü metinler üzerinde kullanılabilir olması için nasıl bir sistem tasarımı gereklidir? Proje bu sorulara hem deneysel model değerlendirmesi hem de çalışan uygulama geliştirme yoluyla yanıt aramıştır.

Çalışmanın ana amacı, çok dilli haber özetleme için kullanılabilecek modelleri eğitmek, değerlendirmek, karşılaştırmak ve bu modelleri pratik bir haber akışı uygulamasına entegre ederek uçtan uca çalışan bir prototip ortaya koymaktır. Bu amaç olabildiğince ölçülebilir biçimde tanımlanmıştır. Model başarısı; ROUGE, BLEU ve METEOR benzeri referans özetle örtüşme metrikleriyle, kaynak metne bağlılık ve tekrar oranı gibi ek davranış metrikleriyle, ayrıca tutarlılık, doğruluk, açıklık, alakalılık ve verimlilik boyutlarını temsil eden birleşik kabiliyet skorlarıyla ölçülmüştür. Uygulama başarısı ise sistemin farklı dillerdeki RSS kaynaklarından haber toplayabilmesi, haber gövdesini çıkarabilmesi, model bazlı özetleri üretebilmesi, bu özetleri veritabanında saklayabilmesi ve kullanıcının dil, konu, ülke, bölge, kaynak, anahtar kelime ve model seçimine göre haberleri görüntüleyebilmesi üzerinden değerlendirilmiştir. Böylece proje yalnızca “iyi özet üreten model” hedefiyle sınırlı kalmamış; modelin gerçek bir kullanıcı arayüzünde ne kadar kullanılabilir hale getirilebildiğini de ölçmeyi hedeflemiştir.

Projenin özgün değeri, model eğitimi, model değerlendirmesi ve uygulama entegrasyonunu tek bir bütün halinde ele almasından gelmektedir. Haber özetleme üzerine yapılan birçok çalışmada model performansı yalnızca veri kümesi üzerindeki otomatik metriklerle incelenirken, bu projede modellerin gerçek haber akışı içinde nasıl davranacağı da dikkate alınmıştır. RSS kaynaklarından gelen haberler, standart ve temiz akademik veri kümelerindeki metinlerden farklıdır; başlık tekrarları, reklam metinleri, tarih ve yazar bilgileri, sosyal medya bağlantıları, eksik paragraflar ve kaynaklara özgü şablon kalıntıları içerebilir. Bu proje, bu tür gürültülü girdileri temizleyerek özetleme modeline daha uygun hale getiren bir işleme hattı kurmuştur. Ayrıca kaynak metin alınamadığında sistemi tamamen durdurmak yerine kontrollü geri dönüş mekanizmaları kullanılmış, ancak düşük kaliteli veya yetersiz içeriklerin sisteme zarar vermesini engellemek için kalite kapıları tasarlanmıştır.

Özgünlük açısından bir diğer önemli nokta, değerlendirmenin tek boyutlu yapılmamasıdır. Haber özetleme modellerinde yalnızca referans özetle kelime örtüşmesine bakmak yanıltıcı olabilir. Bir model referansa yakın kelimeler üretebilir ancak kaynak metinden kopabilir, gereksiz tekrar yapabilir, çok uzun veya çok kısa özetler üretebilir ya da pratik kullanım için yavaş kalabilir. Bu nedenle projede klasik metriklerin yanında kaynak kapsaması, kaynak geri çağırımı, parça kapsaması, tekrar oranı, sıkıştırma oranı ve gecikme süresi gibi ek ölçümler kullanılmıştır. Bu ölçümler daha sonra tutarlılık, doğruluk, açıklık, alakalılık ve verimlilik gibi daha yorumlanabilir başlıklar altında birleştirilmiştir. Böylece modeller yalnızca akademik skorlarına göre değil, gerçek kullanıcı deneyimine etkisi olacak yönleriyle de karşılaştırılmıştır.

Konunun seçilme nedeni, haber tüketiminde zaman, dil ve güvenilirlik problemlerinin aynı anda ortaya çıkmasıdır. Kullanıcılar çoğu zaman birden fazla kaynağı takip etmek, farklı ülkelerden haberleri görmek, yabancı dildeki içerikleri anlamak ve uzun metinleri kısa sürede kavramak ister. Bununla birlikte otomatik özetleme sistemleri, haber gibi hassas bir alanda dikkatli tasarlanmadığında eksik, yanıltıcı veya bağlamdan kopuk çıktılar üretebilir. Bu proje, bu ihtiyacı hem doğal dil işleme hem de yazılım sistemi tasarımı açısından ele alarak anlamlı bir çözüm geliştirmeyi amaçlamıştır. Çok dilli destek, model karşılaştırması ve web uygulaması entegrasyonu birlikte düşünüldüğünde proje, yalnızca deneysel bir model çalışması değil, kullanıcıya dönük bir haber özetleme platformu prototipi niteliği taşımaktadır.

Projenin ana tasarımı, ardışık fakat birbirini besleyen üç katman üzerine kuruludur. İlk katman model katmanıdır. Bu katmanda haber özetleme için uygun transformer tabanlı diziden diziye modeller kullanılmıştır. İngilizce haber özetleme için BART tabanlı modeller, çok dilli haber özetleme için ise mBART ve mT5 tabanlı modeller ele alınmıştır. Proje kapsamında mBART tabanlı iki farklı eğitim hattı oluşturulmuştur. İlk eğitim hattı daha temel ve güvenli bir başlangıç yaklaşımı olarak tasarlanmış, ikinci eğitim hattında ise doğrulama stratejisi, hedef uzunluk, öğrenme oranı, örnekleme yaklaşımı ve en iyi kontrol noktasını seçme gibi kalite odaklı düzenlemeler yapılmıştır. Böylece aynı temel mimari üzerinde farklı eğitim kararlarının model davranışına etkisi karşılaştırılabilir hale getirilmiştir.

İkinci katman değerlendirme katmanıdır. Bu katmanda modeller, çok dilli haber özetleme veri kümesi üzerinde test edilmiştir. Değerlendirmede her modelin farklı dillerdeki çıktıları ayrı ayrı incelenmiş, ardından genel ortalama performansları karşılaştırılmıştır. Klasik metrikler modelin referans özetle ne kadar örtüştüğünü gösterirken, kaynak bağlılığı metrikleri özetin haber metninden ne kadar kopmadığını anlamak için kullanılmıştır. Tekrar oranı ve okunabilirlik sinyalleri özetin kullanıcı tarafından ne kadar rahat okunabileceğine dair fikir vermiştir. Gecikme süresi ve sıkıştırma oranı ise modelin web uygulamasında pratik olarak kullanılabilir olup olmadığını değerlendirmek için önemlidir. Bu yaklaşım, projenin yalnızca model skoru üretmekle kalmayıp, model seçimini uygulama bağlamına göre yapmasını sağlamıştır.

Üçüncü katman uygulama katmanıdır. Web uygulaması Flask tabanlı bir mimariyle geliştirilmiştir. Sistem, farklı diller ve kategoriler için tanımlanmış RSS kaynaklarından haberleri belirli aralıklarla toplar. Toplanan haberlerin bağlantıları üzerinden haber sayfasına gidilir ve mümkünse tam haber gövdesi çıkarılır. Haber gövdesi temizlendikten sonra seçilen özetleme modellerine gönderilir. Üretilen özetler haber kaydıyla birlikte SQLite veritabanına yazılır. Bu tercih, API çağrısı yapıldığında modelin her seferinde yeniden çalıştırılmasını önleyerek arayüzün daha hızlı yanıt vermesini sağlar. Arka planda çalışan zamanlayıcı, haberleri düzenli aralıklarla işler ve yeni özetleri hazır hale getirir. Kullanıcı arayüzü ise hazır özetleri veritabanından okuyarak dil, konu, ülke, bölge, kaynak, anahtar kelime ve model filtresiyle sunar.

Projede kullanılan yöntemler, amaçlarla doğrudan ilişkilidir. Model eğitimi için transformer tabanlı önceden eğitilmiş modellerin seçilmesi, çok dilli özetleme görevinin yüksek veri ve hesaplama gerektiren yapısından kaynaklanmaktadır. Sıfırdan model eğitmek yerine önceden eğitilmiş modelleri haber özetleme görevine uyarlamak, sınırlı donanım ve süre koşullarında daha gerçekçi bir yaklaşımdır. Değerlendirme için çok sayıda metriğin birlikte kullanılması, haber özetleme kalitesinin tek bir skorla temsil edilememesinden kaynaklanmaktadır. Uygulama tarafında RSS tabanlı veri toplama kullanılması ise haberlerin düzenli, güncel ve kaynak bazlı biçimde alınmasına olanak tanır. SQLite kullanımı prototip düzeyinde yeterli, taşınabilir ve kurulumu kolay bir veri saklama çözümü sağlamıştır. Arka plan ingest mimarisi ise model üretim maliyetini kullanıcı isteği anından ayırarak sistemin daha kullanılabilir olmasına katkı vermiştir.

Kaynak ihtiyaçları açısından proje dört ana gruba ayrılabilir. Birinci grup veri kaynaklarıdır. Model eğitimi ve değerlendirmesi için haber özetleme veri kümeleri, uygulama için ise RSS haber kaynakları kullanılmıştır. İkinci grup yazılım kaynaklarıdır. Python, Flask, transformers, PyTorch, RSS ayrıştırma ve haber gövdesi çıkarma kütüphaneleri projenin temel yazılım bileşenlerini oluşturmuştur. Üçüncü grup donanım kaynaklarıdır. Transformer tabanlı modellerin eğitimi ve değerlendirilmesi GPU gereksinimi doğurmuştur; buna karşılık uygulama tarafında hazır modellerin çalıştırılması için daha sınırlı ama yine de dikkatli yönetilmesi gereken işlemci, bellek ve depolama kaynakları gerekmektedir. Dördüncü grup ise sistem kaynaklarıdır. Haberlerin saklanması için veritabanı, model dosyalarının tutulması için yerel depolama, arka plan işlerinin izlenmesi için loglama ve düzenli ingest için zamanlayıcı altyapısı kullanılmıştır.

Projeden elde edilen sonuçlar, model seçiminin tek bir metrik üzerinden yapılamayacağını göstermiştir. Hazır mT5 tabanlı model genel değerlendirmede güçlü bir karşılaştırma noktası oluşturmuş ve birçok klasik metrikte yüksek sonuçlar vermiştir. Proje kapsamında eğitilen iki mBART modeli ise birbirine yakın genel performans göstermiş, ancak davranış profilleri farklılaşmıştır. İlk mBART modeli kaynak metne daha bağlı ve daha güvenli özetler üretme eğilimindeyken, ikinci mBART modeli bazı klasik örtüşme metriklerinde ve hız tarafında küçük avantajlar sağlamıştır. Bu sonuç, haber özetleme sistemlerinde “en iyi model” kararının kullanım amacına bağlı olduğunu ortaya koymaktadır. Eğer öncelik kaynak sadakati ve daha düşük olgusal risk ise daha kaynak-bağlı model tercih edilebilir. Eğer öncelik daha hızlı, daha kısa ve daha serbest yeniden ifade edilen özetler ise diğer model daha uygun olabilir.

Uygulama entegrasyonu açısından proje, model çıktılarının gerçek haber akışı üzerinde kullanılabilir hale getirilebildiğini göstermiştir. Sistem, farklı dillerdeki kaynakları okuyabilmekte, haber metinlerini çıkarabilmekte, özetleri model bazlı olarak önceden hesaplayabilmekte ve kullanıcıya filtrelenebilir bir arayüzde sunabilmektedir. Ayrıca çeviri desteği, görsel gösterimi, kaynak bağlantısı, anahtar kelime araması, konu ve bölge filtreleri gibi özellikler uygulamayı yalnızca teknik bir demo olmaktan çıkarıp daha işlevsel bir haber takip aracına yaklaştırmıştır. Arka planda çalışan ingest yapısı sayesinde haber toplama ve özet üretme işlemleri düzenli biçimde yapılmakta, API ise kullanıcıya hazır sonuçları döndürmektedir. Bu mimari, özetleme modellerinin yüksek çalışma maliyetini azaltmak ve kullanıcı deneyimini iyileştirmek açısından önemlidir.

Genel olarak proje, çok dilli haber özetleme problemini model, değerlendirme ve uygulama eksenlerinde ele almıştır. Çalışma sonucunda farklı özetleme modellerinin güçlü ve zayıf yönleri ölçülmüş, bu modellerin gerçek haber kaynaklarıyla çalışan bir sistem içinde nasıl kullanılabileceği gösterilmiş ve kullanıcıya dönük bir web uygulaması prototipi geliştirilmiştir. Projenin en önemli çıktısı, haber özetleme kalitesinin yalnızca referans metin benzerliğiyle değil, kaynak sadakati, okunabilirlik, tekrar oranı, hız ve sistem entegrasyonu gibi çok boyutlu ölçütlerle değerlendirilmesi gerektiğini göstermesidir. Bu yönüyle proje, hem doğal dil işleme alanında deneysel bir çalışma hem de gerçek kullanım senaryosuna yaklaşan bir yazılım sistemi geliştirme çalışması olarak değerlendirilebilir.